\documentclass[11pt]{article}
% === core_packages.tex ===
\usepackage{amsmath, amssymb, amsthm, amsfonts}
\usepackage{geometry}
\usepackage{hyperref}
\usepackage{graphicx}
\usepackage{physics}
\usepackage{booktabs}
\usepackage{listings}
\usepackage{bookmark}
\usepackage{amscd}
\usepackage{tikz-cd}
\usepackage{mathtools}
\usepackage{xspace}
\usepackage[backend=biber, style=numeric]{biblatex}
\usepackage{tcolorbox}
\tcbuselibrary{theorems}
\usepackage{parskip}
\usepackage{tikz}
\usepackage{etoolbox}
\usetikzlibrary{arrows.meta, shapes, positioning, calc, cd, shapes.geometric, backgrounds}
\usepackage[en-US]{datetime2}
\usepackage{array}


% --- Math & Ontology Macros ---
\newcommand{\entropy}{S}
\newcommand{\opponent}{opponent}
\newcommand{\aside}[1]{[\emph{\opponent}: ``#1'']}
\newcommand{\paper}[1]{[\emph{#1}: ``#1'']}
\newcommand{\iame}{Real\textsuperscript{e}}
\newcommand{\iameverything}{Real\textsuperscript{everything}}
\newcommand{\fabric}{voxel}
\newcommand{\Fabric}{Voxel}
\newcommand{\pdflink}[1]{}



% --- Global Layout Defaults ---
\geometry{letterpaper, margin=1in}
\hypersetup{
    colorlinks=true,
    linkcolor=black,
    citecolor=black,
    filecolor=magenta,
    urlcolor=blue
}
\hfuzz=5pt \vfuzz=5pt \hbadness=10000 \vbadness=10000

\author{\iame}


\usepackage{tocloft}

% --- Article Specific Theorems ---
\newtheorem{lemma}{Lemma}
\theoremstyle{definition}
\newtheorem{axiom}{Axiom}
\newtheorem{theorem}{Theorem}
\newtheorem{corollary}{Corollary}
\newtheorem{definition}{Definition}
\newtheorem{proposition}{Proposition}
\newtheorem{conjecture}{Conjecture}


\makeatletter
\let\remark\@undefined
\let\endremark\@undefined
\makeatother
\theoremstyle{remark}
\newtheorem*{remark}{Remark}

\newtcbtheorem[number within=section]{principle}{Principle}%
{colback=blue!5,colframe=blue!50!black,fonttitle=\bfseries}{principle}

% --- Article TOC Layout ---
\setlength{\cftbeforesecskip}{1.0em}
\renewcommand{\cftsecleader}{\cftdotfill{\cftdotsep}}
\renewcommand{\cftsubsecleader}{\cftdotfill{\cftdotsep}}


\addbibresource{../references.bib}

\title{Resolution of Spacetime}

\author{Juha Meskanen}
\date{June 2026}

\begin{document}

\maketitle

\begin{abstract}
A static, finite informational universe is defined by \(n\) bits.  
The \(2^n\) arrangements of those bits constitute the entire 
configuration space. Time is not fundamental; it emerges as the 
ordinal index of a random walk through the most compressible 
configurations. Because the total information is finite, both 
spatial and temporal resolution of any emergent spacetime must 
likewise be finite.  

There is no external metaphysics and no predefined rule that 
selects one resolution convention over another. The only selection 
principle available is typicality: the measure-dominated, most 
probable outcome under the natural dynamics on the configuration 
space. Spectral complexity---the informational cost of encoding 
the modes of a wavefunction---supplies the compressibility measure 
that identifies the typical path.  

Under that path the elementary bits behave as structureless spacetime 
fabric. Emergent microstructures appear with a characteristic 
hump-shaped probability, consume part of the finite bit budget, and 
thereby reduce the observable resolution. The resulting geometry is 
a chain whose links are bits and whose knots are the microstructures; 
each knot shortens the chain. The same finite-budget logic applied 
to clocks yields differential aging. The typical evolution of such a 
chain produces a three-stage expansion profile without fine-tuning.
\end{abstract}

\section{The finite informational premise}

We assume a static, timeless universe whose entire content is a 
finite collection of \(n\) bits. Bits themselves are not 
ontologically privileged; they are merely one convenient binary 
representation of a finite amount of information. Any other finite 
alphabet would serve equally well.

The configuration space therefore consists of exactly \(2^n\) 
distinct arrangements. No continuous manifold, no continuous time 
parameter, and no infinite reservoir of information are postulated. 
All that exists is the set of \(2^n\) discrete states together with 
the combinatorial structure that relates them.

\section{Emergence of ordinal time}

Time is the ordinal index of a random walk through the most 
compressible configurations. The minimal---and hence most 
compressible---path is the path of least change between consecutive 
configurations. That path is realised by a fair random bit-flip 
mutation of an \(n\)-bit string: at each discrete tick one of the 
\(n\) bits is chosen uniformly and flipped.

This dynamics is the Ehrenfest process. Let \(m\) be the number of 
1-bits and \(p=m/n\). The expected change in one tick is
\begin{equation}
  \mathbb{E}[\Delta m]=\frac{n-m}{n}-\frac{m}{n}=1-2p.
\end{equation}
In the continuum limit one obtains the ordinary differential equation
\begin{equation}
  \frac{dp}{dt}=\frac{1-2p}{n},
\end{equation}
where \(t\) is the raw flip count (``bit-flip time''). Starting from 
the all-zero state \(p(0)=0\) the solution is
\begin{equation}
  p(t)=\frac12\bigl(1-e^{-2t/n}\bigr).
\end{equation}
The relaxation timescale is therefore \(O(n)\) ticks. Rescaling to 
the dimensionless coordinate \(\tau=t/n\) yields the universal curve
\begin{equation}
  p(\tau)=\frac12\bigl(1-e^{-2\tau}\bigr),\qquad\tau\in[0,\infty).
\end{equation}
The associated Shannon entropy per bit is
\begin{equation}
  \frac{S(\tau)}{n}=H_2\bigl(p(\tau)\bigr),
\end{equation}
where \(H_2(q)=-q\log_2 q-(1-q)\log_2(1-q)\).

A further change of variable
\begin{equation}
  x(\tau)=1-e^{-2\tau}\in[0,1)
\end{equation}
produces a bounded time coordinate whose entropy is simply
\begin{equation}
  \frac{S(x)}{n}=H_2\bigl(x/2\bigr).
\end{equation}
This is the universal entropy curve that any physical time scale is 
ultimately mapped onto.

Compressibility is measured by spectral complexity: the number of 
bits required to encode the modes (frequencies and phases) of a 
wavefunction. The Ehrenfest trajectory is the path of least spectral 
complexity; it is therefore the typical path under the natural measure 
on configuration space. Solomonoff induction, when formulated with 
spectral rather than Kolmogorov complexity, preferentially samples 
exactly this trajectory.

\section{Resolution of spacetime}

Because only \(n\) bits of information exist, any emergent geometry 
can possess only finitely many distinguishable spatial cells and only 
finitely many distinguishable temporal ticks. This is the principle 
of \emph{resolution of spacetime}:

\begin{quote}
  The spatial resolution of any region and the temporal resolution 
  of any clock are bounded by the finite informational budget of 
  the underlying bit-string.
\end{quote}

No external rule dictates how that budget is to be partitioned into 
space and time. The only admissible answer is typicality: the 
partition that dominates the measure under the Ehrenfest dynamics.

\subsection{Maximal theoretical resolution}

In the absence of any emergent structure the entire entropy budget 
is available for geometry. The maximal theoretical resolution is 
therefore the pure entropy curve itself:
\begin{equation}
  R_{\mathrm{max}}(\tau)=\frac{S(\tau)}{n}=H_2\bigl(p(\tau)\bigr).
\end{equation}
This is the resolution an internal observer would measure if the 
bit-string remained completely unstructured.

\subsection{Spatial resolution -- the chain of knots}

Bits may be pictured as structureless spacetime fabric---the links of 
a chain. Microstructures form when local windows of the chain adopt 
particular compositions. Each microstructure is a knot: it is made 
from the same links that previously contributed to the length of the 
chain, and once the knot is tied those links are no longer available 
to resolve distance. If a knot consumes \(w\) bits, the observable 
spatial resolution drops by \(w\) for every such knot. The residual 
budget
\begin{equation}
  R_Q(\tau)=R_{\mathrm{max}}(\tau)-\frac{1}{n}\sum\text{(bits locked in knots)}
\end{equation}
is the only quantity still available to an internal observer as 
spatial resolution.

\subsection{Temporal resolution -- differential aging}

A structure of width \(w\) can advance its own proper time only after 
a full set of \(w\) elementary flips has occurred. Its retarded clock 
is therefore the sample-and-hold quantisation
\begin{equation}
  \tau_{\mathrm{local}}(w)=w\cdot\bigl\lfloor\tau/w\bigr\rfloor,
\end{equation}
with associated lapse \(1/w\). Heavier knots age more slowly with 
respect to the universal raw coordinate \(\tau\). The informational 
backlog that the raw clock has already registered but the retarded 
clock has not yet seen is the pending entropy of that scale. Pending 
bits remain part of the finite budget; they simply have not yet been 
experienced by the local observer.

The same chain-of-knots analogy applies temporally: each knot 
introduces a coarser ticking rate, stretching the perceived duration 
of events that occur at that scale.

\section{Emergence of patterns and the dominance of the hump}

For a fixed window of length \(L\) with composition \((a,b)\) 
(\(a+b=L\)) the probability of that exact pattern under the 
independent-bit continuum limit is
\begin{equation}
  P(\text{pattern}\mid\tau)=p(\tau)^a\bigl(1-p(\tau)\bigr)^b.
\end{equation}
Three exact results follow at once.

\begin{enumerate}
  \item Zero/one-start complementation symmetry is exact for all \(\tau\).
  \item Only the composition \((a,b)\) matters; every permutation of 
        a given composition yields the identical curve.
  \item On the reachable domain \(p\in[0,1/2]\) there are exactly 
        three qualitative shapes, classified by the sign of \(a-b\):
        \begin{itemize}
          \item \(a>b\): monotonic rise (peak lies beyond equilibrium);
          \item \(a=b\): monotonic rise that flattens at the boundary;
          \item \(a<b\): genuine interior hump.
        \end{itemize}
\end{enumerate}

Only the hump class (\(a<b\)) produces a rise-peak-decline profile. 
Moreover, when each composition is weighted by its combinatorial 
multiplicity \(\binom{L}{a}\), the hump class dominates the measure 
for almost the entire relaxation. The window occupancy is binomial 
with mean \(Lp(\tau)<L/2\) at every finite \(\tau\); consequently the 
mass of the hump region exceeds the mass of the non-hump region until 
a crossover \(\tau_c\) that recedes toward infinity as \(L\) grows. 
For any practically observable window length the hump is therefore 
the typical outcome.

The peak location, amplitude and asymptotic decay of a hump are fixed 
by elementary formulae:
\begin{align}
  p^*&=\frac{a}{a+b},&
  \tau^*&=-\frac12\ln(1-2p^*),\\
  \frac{\text{peak}}{\text{plateau}}&=2^{L\bigl(1-H_2(p^*)\bigr)}.
\end{align}
No free parameters are required.

\section{The typical expansion profile}

Because microstructures consume part of the finite bit budget, the 
observable resolution \(R_Q\) is necessarily smaller than the maximal 
resolution \(R_{\mathrm{max}}\). The typical history of a chain of 
knots therefore exhibits three successive regimes:

\begin{itemize}
  \item rapid early growth of \(R_Q\) while knots are still rare;
  \item matter-loading slowdown once the hump populations lock a 
        non-negligible fraction of the budget into knots;
  \item late recovery as the hump populations decline and the residual 
        informational floor is approached.
\end{itemize}

This three-stage profile is not imposed by hand. It is the measure-
dominated behaviour of the only dynamics that is typical under spectral 
complexity. The fine-tuning problem is thereby dissolved: the observed 
expansion history is simply what a typical finite informational universe 
looks like from the inside.

\section{Conclusion}

Spacetime resolution is the direct geometric consequence of a finite 
informational substrate. Once \(n\) is finite, both the number of 
distinguishable spatial cells and the number of distinguishable 
temporal ticks are bounded. The pure entropy curve supplies the maximal 
theoretical resolution; every emergent knot reduces that resolution by 
the number of bits it consumes. Temporal resolution is coarsened by the 
same knots through retarded clocks whose lapse is the reciprocal of 
their width.

All concrete choices---the Ehrenfest path, the dominance of the hump, 
the conversion of residual budget into spatial length, the differential 
aging of clocks---are selected solely by typicality. There is no 
external rulebook. The only answer to any question about how spacetime 
is resolved is: the answer that dominates the measure.

Continuous geometries appear only in the formal limit \(n\to\infty\). 
For any realistic finite \(n\) the resolution of spacetime remains 
finite, discrete, and fully determined by the underlying bit-string 
and the typical path through its configurations.

\printbibliography
\end{document}
